\clearpage
\SourceLeaf{411}%
\fancyhead[L]{\small VERSCHIEDENE AUFSÄTZE.}
\pdfbookmark[1]{Lösung einer geometrischen Aufgabe.}{aufloesung-geometrische-1823}
\begin{center}
\rule{.8\textwidth}{.4pt}\\[4pt]
Astronomische Nachrichten. Nr. 42. November 1823.\\[4pt]
\rule{.8\textwidth}{.4pt}\\[12pt]
\textit{Lösung einer geometrischen Aufgabe.}
\end{center}
\hspace*{\parindent}\{In der Beschreibung der Leipziger Sternwarte von Professor \textsc{Möbius} findet sich auf S.\,61 die folgende kleine geometrische Aufgabe.
\begin{quote}
Beliebige 5 Punkte \(A,B,C,D,E\) in einer Ebene werden paarweise durch Geraden verbunden. Man kennt die Flächeninhalte der so entstehenden 5 Dreiecke \(EAB,ABC,BCD,CDE,DEA\) und soll daraus den Flächeninhalt des Fünfecks \(ABCDE\) bestimmen.
\end{quote}
Hofrat \textsc{Gauss}, der diesen Sommer einige Wochen bei mir verbrachte, schrieb, als er das Buch sah, die folgende Lösung hinein.

Man bezeichne die 5 Punkte mit \(1,2\ 3,4,5\), die
\begin{center}
\begin{tabular}{lrl}
Winkel&\(213\)&mit \(p\)\\
&\(214\)&— \(q\)\\
&\(215\)&— \(r\)\\\hline
Seiten&\(12\)&— \(t\)\\
&\(13\)&— \(u\)\\
&\(14\)&— \(v\)\\
&\(15\)&— \(w\)\\\hline
Dreiecke&\(123\)&— \(a\)\\
&\(234\)&— \(b\)\\
&\(345\)&— \(c\)\\
&\(451\)&— \(d\)\\
&\(512\)&— \(e\)\\
&\(124\)&— \(x\)\\
&\(134\)&— \(y\)\\
&\(135\)&— \(z\)\\\hline
Fünfeck&\(12345\)&— \(\omega\)
\end{tabular}
\end{center}
